% Einheit 1 – Proportionale Funktion (Nr. 1–8)
\setcounter{aufgabe}{0}
\einheitenkopf[e1][1 Proportionale Funktion]{Einheit 1 von 5 · Proportionale Funktion}
\zweigzeile{Hier lernst du, Funktionen $y = m \cdot x$ zu zeichnen, abzulesen und aufzustellen · neu in diesem Jahr · P10 oft · baut auf: Dreisatz, proportionale Zuordnung, Koordinatensystem}

\verfahren{Wertetabelle und Graph von $y = m \cdot x$}

\begin{aufgabe}{Ich kann zu einer Gleichung $y = m \cdot x$ eine Wertetabelle ausfüllen.\newline Fülle die Wertetabelle aus.}
\begin{teile}
\teil $y = 2x$ \quad \wertetabelle{x}{y}{0,1,2,3}
\teil $y = 3x$ \quad \wertetabelle{x}{y}{0,1,2,3}
\teil $y = x$ \quad \wertetabelle{x}{y}{0,1,2,3}
\teil $y = 5x$ \quad \wertetabelle{x}{y}{0,1,2,3}
\teil $y = 1{,}5x$ \quad \wertetabelle{x}{y}{0,1,2,3}
\teil $y = \frac{1}{2}x$ \quad \wertetabelle{x}{y}{0,1,2,3}
\teil $y = -2x$ \quad \wertetabelle{x}{y}{0,1,2,3}
\end{teile}
\end{aufgabe}

\begin{aufgabe}{Ich kann die Gerade zu $y = m \cdot x$ zeichnen.\newline Zeichne die Gerade in das Koordinatensystem. Deine Tabellen aus Nr.~1 helfen.}
\begin{teile}
\teil $y = 2x$
\teil $y = \frac{1}{2}x$
\teil $y = -2x$
\end{teile}

\begin{ksys}[xmin=-3,xmax=3,ymin=-6,ymax=6]
\end{ksys}

\end{aufgabe}

\verfahren{Werte ablesen und berechnen}

\begin{aufgabe}{Ich kann Werte am Graphen einer proportionalen Funktion ablesen und berechnen.\newline Die Gerade $g$ gehört zu $y = 1{,}5x$. Lies am Graphen ab.}

\begin{ksys}[xmin=-5,xmax=5,ymin=-6,ymax=6,ablesen]
\gerade[3.2]{1.5}{0}{g}
\end{ksys}

\begin{teilezwei}
\tz $y$ bei $x = 2$: \quad \feld{y} & \tz $y$ bei $x = 4$: \quad \feld{y} \\
\tz $y$ bei $x = -2$: \quad \feld{y} & \tz $x$ bei $y = -6$: \quad \feld{x} \\
\anweisung{Rechne mit der Gleichung. Diese Werte liegen außerhalb der Zeichnung.}
\tz $y$ bei $x = 10$: \quad \feld{y} & \tz $x$ bei $y = 21$: \quad \feld{x} \\
\end{teilezwei}
\end{aufgabe}

\verfahren{Proportionale Funktion erkennen}

\begin{aufgabe}{Ich kann entscheiden, ob eine Tabelle oder ein Text zu einer proportionalen Funktion gehört.\newline Kreuze an. Wenn ja, gib die Gleichung an.}
\begin{teile}
\teil \wertetabelle[3,6,9,12]{x}{y}{1,2,3,4} \quad \janein \quad $y = $ \leerfeld
\teil \wertetabelle[2,4,7,8]{x}{y}{1,2,3,4} \quad \janein \quad $y = $ \leerfeld
\teil \wertetabelle[5,10,15,20]{x}{y}{2,4,6,8} \quad \janein \quad $y = $ \leerfeld
\teil \wertetabelle[5,6,7,8]{x}{y}{1,2,3,4} \quad \janein \quad $y = $ \leerfeld
\end{teile}
\end{aufgabe}

\begin{aufgabe}{Ich kann entscheiden, ob eine Tabelle oder ein Text zu einer proportionalen Funktion gehört – weiter.}
\begin{teile}
\teil \wertetabelle[-1,-2,-3,-5]{x}{y}{2,4,6,10} \quad \janein \quad $y = $ \leerfeld
\teil Ein Kilogramm Kirschen kostet 4\,€. $x$: Masse in kg, $y$: Preis in €. \newline \janein \quad $y = $ \leerfeld
\teil Ein Handytarif kostet 5\,€ im Monat und 0,10\,€ je Minute. $x$: Minuten, $y$: Preis in €. \newline \janein \quad $y = $ \leerfeld
\end{teile}
\end{aufgabe}

\verfahren{Werte mit dem Dreisatz kontrollieren}

\begin{aufgabe}{Ich kann einen Funktionswert mit dem Dreisatz kontrollieren.\newline $x$ ist die Masse in kg, $y$ der Preis in €. Berechne $y$ mit der Gleichung und kontrolliere mit dem Dreisatz.}
\rechenplatz{3}

\noindent\begin{minipage}[t]{0.32\linewidth}
\begin{teile}\teil $y = 4x$, \ $x = 3$\end{teile}
\begin{dreisatz}{kg}{€}
\dsz{1}{\dsleer}
\dsp{\cdot 3}{\cdot 3}
\dsz{3}{\dsleer}
\end{dreisatz}
\end{minipage}\hfill
\begin{minipage}[t]{0.32\linewidth}
\begin{teile}\teil $y = 2{,}5x$, \ $x = 4$\end{teile}
\begin{dreisatz}{kg}{€}
\dsz{1}{\dsleer}
\dsp{\cdot 4}{\cdot 4}
\dsz{4}{\dsleer}
\end{dreisatz}
\end{minipage}\hfill
\begin{minipage}[t]{0.32\linewidth}
\begin{teile}\teil $y = 6x$, \ $x = 1{,}5$\end{teile}
\begin{dreisatz}{kg}{€}
\dsz{1}{\dsleer}
\dsp{\cdot 1{,}5}{\cdot 1{,}5}
\dsz{1{,}5}{\dsleer}
\end{dreisatz}
\end{minipage}
\end{aufgabe}

\verfahren{Prüfen, begründen, anwenden}

\begin{aufgabe}{Ich finde den Fehler beim Bestimmen des Faktors.\newline Jede Tabelle gehört zu einer proportionalen Funktion. Finde den Fehler, benenne ihn und korrigiere die Rechnung.}
\begin{teile}
\teil \wertetabelle[7,14,21]{x}{y}{2,4,6} \quad Max rechnet:
\rechnung{m &= 14 - 7 = 7 \\ y &= 7x}
\teil \wertetabelle[-6,-12,-18]{x}{y}{3,6,9} \quad Ben rechnet:
\rechnung{m &= -12 - (-6) = -6 \\ y &= -6x}
\end{teile}
\end{aufgabe}

\begin{aufgabe}{Ich kann begründen, ob eine Gerade zu einer proportionalen Funktion gehört.\newline Entscheide ohne Rechnung und begründe.}
\begin{teile}
\teil Eine Gerade geht durch den Punkt $(0 \,|\, 3)$. Kann sie zu einer Gleichung $y = m \cdot x$ gehören?
\teil Tim sagt: „Bei $y = 4x$ verdoppelt sich $y$, wenn ich $x$ verdopple.“ Hat Tim recht?
\end{teile}
\end{aufgabe}

\begin{aufgabe}{Ich kann eine Sachaufgabe zur proportionalen Funktion lösen.\newline Ein leeres Aquarium wird mit einem Schlauch gefüllt. In jeder Minute fließen 15 Liter hinein. $x$ ist die Zeit in Minuten, $y$ die Wassermenge in Litern.}
\begin{teile}
\teil Stelle die Gleichung auf. \quad $y = $ \leerfeld
\teil Zeichne den Graphen für die ersten 8 Minuten.
\teil Um wie viel Liter wächst $y$, wenn $x$ um 1 größer wird? Wo siehst du das im Graphen?
\end{teile}

\begin{ksys}[xmin=0,xmax=8,ymin=0,ymax=120,ystep=15,xlabel=$x$ in min,ylabel=$y$ in l]
\end{ksys}

\end{aufgabe}
